% Kanonische Aussagen und Referenzziele stammen ausschließlich aus Band 19.
\begin{luacode*}
  local directory = "registry/cuts/"
  thmlookup.registry_path = directory .. tex.jobname .. ".registry.tsv"
  thmlookup.debug_path = directory .. tex.jobname .. ".debug.log"
  thmlookup.prepare_run()
  for _, predecessor in ipairs(banddeps.graph.B19.predecessors) do
    local base = banddeps.graph[predecessor].artifact_base
    tex.sprint("\\externaldocument{" .. base .. "}[../_" .. predecessor .. ".pdf]")
  end
  tex.sprint("\\externaldocument{registry/_B19}[../_B19.pdf]")
  banddeps.load_predecessor_registries("B19")
  thmlookup.load_registry_file("registry/_B19.registry.tsv")
\end{luacode*}
\renewcommand{\FormulaBandID}{19}
\setcounter{file}{19}
\providecommand{\CutsCurrentNumber}{19.0.0.0}
% Split-Tabellen setzen ihren Elternsatz aus diesem Standardzählerbild.
\renewcommand{\theformulaThm}{\CutsCurrentNumber}
\renewcommand{\theHtable}{cuts.\arabic{chapter}.\arabic{table}}
\renewcommand{\theHfigure}{cuts.\arabic{chapter}.\arabic{figure}}
\renewcommand{\theHproofpartnr}{cuts.\mxProofPartParent.\arabic{proofpartnr}}
\providecommand{\RealCut}[1]{\operatorname{Cut}_{\mathbb Q}\!\left(#1\right)}
\providecommand{\RealHat}[1]{\widehat{#1}}
\providecommand{\RealEmb}{\iota_{\mathbb R}}
\providecommand{\RealLe}{\mathrel{\leq_{\mathbb R}}}
\providecommand{\RealLt}{\mathrel{<_{\mathbb R}}}

% Derselbe Quelltext wird hier als importierte Aussage ausgegeben. Keine
% zusätzliche Deklaration, kein eigener Formelzähler, kein neues Registryziel.
\newcommand{\CutsCanonicalStatement}[1]{%
  \par\Needspace{12\baselineskip}%
  \hypertarget{cuts.proof.#1}{}%
  \directlua{
    local identity = thmlookup.by_id["\luaescapestring{#1}"]
    local record = identity and thmlookup.by_label[identity.label]
    if not record then
      tex.error("Missing canonical cut statement: \luaescapestring{#1}")
    else
      tex.sprint("\\gdef\\CutsCurrentNumber{" .. record.number .. "}")
    end
  }%
  \ProofPartSetParent{\CutsCurrentNumber}%
  \setcounter{proofpartnr}{0}%
  \RegistryStatement[Aussage aus Band 19]{#1}%
}
\RenewDocumentCommand{\FormulaDefDeltaK}{o m m +m}{\CutsCanonicalStatement{#3}}
\RenewDocumentCommand{\FormulaAxiomDeltaK}{o m m +m}{\CutsCanonicalStatement{#3}}
\RenewDocumentCommand{\FormulaThmDeltaK}{o m m +m}{\CutsCanonicalStatement{#3}}

% Die H-Aussagen stehen ebenso ausschließlich im Bandregister. Die shared
% Tabellenkörper benutzen ihre kanonischen Nummern, setzen aber keine Labels.
\RenewDocumentCommand{\mxProofPartOpen}{m o m o o}{%
  \par\Needspace{9\baselineskip}\medskip
  \stepcounter{proofpartnr}%
  \edef\mxPPLbl{thm:pp:\mxProofPartParent:\arabic{proofpartnr}}%
  \noindent
  \IfNoValueTF{#2}{\textbf{Teil}}{\textbf{#2}}%
  \hfill\ProofReferencesWithoutRecording{\ThmLookupEmitRef{theorem}{\mxPPLbl}}%
  \par\[#3\]\par\small
  \setcounter{proofstepnr}{0}%
  \setlength{\tabcolsep}{2pt}%
  \renewcommand{\arraystretch}{1.15}%
  \begin{longtable}{#1}%
}

% Auch der allgemeine Beweisbeginn und unnummerierte Teilüberschriften
% bleiben mit wenigstens einigen folgenden Tabellenzeilen zusammen.
\AddToHook{env/tabproofsplitwide/before}{\par\Needspace{13\baselineskip}}
\AddToHook{env/tabproofsplit/before}{\par\Needspace{13\baselineskip}}
\let\CutsOriginalProofPartWide\proofpartwide
\renewcommand{\proofpartwide}[1]{%
  \par\Needspace{6\baselineskip}\CutsOriginalProofPartWide{#1}}
